\section{Read the measured tradeoff}
The table records the checked scalar baseline and the checked tiled prototype for
sixteen rows, sixteen reduction positions and sixteen output columns. Inputs contain
small signed integers; every output is verified before its measurements are accepted.
The cache begins cold and uses the aligned stack and whole-function boundary described
above. Repeating the baseline with different descriptor tile values leaves its work
unchanged, because that implementation does not use \(T\).

\begin{table}[htbp]
  \centering
  \caption{Whole-call work for the sixteen-by-sixteen product in the current prototypes.}
  \tagpdfsetup{table/header-rows={1}}
  \small
  \begin{tabular}{lrrrrr}
    \toprule
    Method   & Instructions & Notices & Full misses & Direct & Two-way \\
    \midrule
    Baseline & 41571        & 8744    & 5388        & 4880   & 4620    \\
    Tile 1   & 171703       & 25128   & 5388        & 5644   & 4828    \\
    Tile 2   & 86447        & 16936   & 1804        & 3004   & 2276    \\
    Tile 3   & 72036        & 14888   & 1855        & 3701   & 3597    \\
    Tile 4   & 59419        & 12840   & 1068        & 2324   & 2544    \\
    Tile 8   & 48293        & 10792   & 1548        & 5300   & 5004    \\
    \bottomrule
  \end{tabular}
\end{table}

Tile four has the fewest fully associative misses among these measured tile edges.
It executes more instructions than the baseline, including loop control, repeated
address formation and output partial-sum work. Tile eight executes fewer instructions
than tile four but has more fully associative misses and more direct-mapped misses
than the baseline. A single ranking cannot represent all three observations.

The fully associative cache is not guaranteed to have fewer misses than every
other placement for an arbitrary trace under LRU. Direct mapping can happen to
retain a line in one set while a global LRU decision evicts it. The table's baseline
is such a trace. Associativity changes the replacement choices as well as the set
constraints; compare measured traces rather than assuming an ordering for all inputs.

No host-time or cycle-speedup claim follows from this table. These measurements
count retired simulator instructions and cache-model events. A real processor's
execution time would depend on its instruction implementation, scheduling and memory
system. Host interpreter time depends on the simulator's implementation and measurement
setup. Measure either quantity separately before reporting an elapsed-time result.

\begin{MARSExercise}{Change one experimental variable}
  Repeat the comparison for \((M,K,N)=(15,13,11)\), keeping the allocation bases,
  aligned stack and cache options fixed. Include tile three and tile four. Explain
  which tiles clip, then compare notices and misses separately. Finally move only the
  right input by sixteen bytes and identify the table columns that can change.
\end{MARSExercise}
